\documentclass[12pt,a4paper]{article}
\usepackage{amsmath,amssymb,amsthm,amsfonts}
\usepackage{geometry}
\usepackage{hyperref}
\usepackage{fancyhdr}
\usepackage{titlesec}
\usepackage{enumitem}

\geometry{margin=2.5cm}
\pagestyle{fancy}
\fancyhf{}
\fancyhead[L]{\small Generalized Helix for Automorphic \(L\)-functions}
\fancyhead[R]{\small Hilbert--P\'olya Blueprint}
\fancyfoot[C]{\thepage}
\renewcommand{\headrulewidth}{0.4pt}

\title{\textbf{Generalized Riemann Helix and Dirac Operators\\[0.3em]
for Automorphic \(L\)-functions}\\[0.5em]
\large A Geometric Hilbert--P\'olya Blueprint}
\author{}
\date{\today}

\begin{document}
\maketitle

\begin{abstract}
We extend the construction of the compactified Riemann helix --- a cylindrical curve whose torsion encodes the prime numbers via the explicit formula for the Riemann zeta function --- to a general automorphic \(L\)-function. For any \(L\)-function with a functional equation and an explicit formula, we define a smooth phase function \(\phi_L(t)\) that locks non‑trivial zeros onto a single generator and carries the arithmetic coefficients in its torsion measure. After a logarithmic compactification we obtain a Dirac operator with point interactions located at the logarithms of the arithmetic conductors, whose regularised spectral determinant is conjectured to equal the completed \(L\)-function. The eigenvalues of this operator would then be the imaginary parts of the zeros, giving a uniform geometric realisation of the Hilbert--P\'olya programme for all principal \(L\)-functions.
\end{abstract}

\section{Introduction}

The Hilbert--P\'olya conjecture asserts the existence of a self‑adjoint operator whose eigenvalues coincide with the non‑trivial zeros of the Riemann zeta function. In recent work, Connes, Consani and Moscovici~\cite{CCM2025} have constructed a spectral triple based on the scaling operator on the adèle class space, whose truncated Weil quadratic form is conjectured to encode the Riemann zeros in its ground‑state eigenvector. Independently, the author has observed that the same data can be geometrised as a \emph{Riemann helix} --- a smooth cylindrical curve with a specially tuned torsion --- whose classical closed geodesics are the prime logarithms and whose quantised geodesic flow yields a Dirac operator with point interactions at those logarithms.

The purpose of this note is to show that the entire geometric apparatus generalises verbatim to any ``reasonable'' \(L\)-function possessing a functional equation and an explicit formula. The only required input is the sequence of non‑trivial zeros and the arithmetic coefficients appearing in the Euler product. The resulting \emph{generalised Riemann helix} provides a uniform candidate Hilbert--P\'olya operator for Dirichlet \(L\)-functions, modular form \(L\)-functions, and more generally for automorphic \(L\)-functions on \(\mathrm{GL}_n\).

\section{Data of an \(L\)-function}

Let \(L(s)\) be an \(L\)-function of degree \(d\), with a completed \(L\)-function
\[
\Lambda(s) = N^{s/2} \prod_{j=1}^d \Gamma_{\mathbb{R}}(s+\mu_j) \; L(s),
\qquad \Gamma_{\mathbb{R}}(s) = \pi^{-s/2}\Gamma(s/2),
\]
satisfying the functional equation \(\Lambda(s) = \varepsilon \Lambda(1-\overline{s})\) with \(|\varepsilon|=1\). We assume that all non‑trivial zeros \(\rho = \tfrac12 + i\gamma\) lie on the critical line (GRH), although the geometric construction does not strictly require it --- off‑line zeros would translate into a complex curve.

The explicit formula (Weil, Barner) relates a sum over zeros to a sum over prime powers:
\[
\sum_{\rho} h(\rho) = \sum_{p^m} \frac{c(p^m)}{p^{m/2}}\,h(m\log p) + \text{archimedean terms},
\tag{1}
\]
where \(c(p^m)\) are the Dirichlet coefficients of the logarithmic derivative \(-\frac{L'}{L}(s)\). For example, for \(\zeta(s)\) we have \(c(p^m) = \log p\) (von Mangoldt function \(\Lambda(p^m)\)); for an elliptic curve \(L\)-function, \(c(p)\) are related to the trace of Frobenius.

\section{The Generalised Riemann Helix}

Let \(0 < \gamma_1 < \gamma_2 < \cdots\) be the positive imaginary parts of the zeros of \(L(s)\). Choose a smooth, strictly increasing function \(\phi_L:[\gamma_1,\infty)\to\mathbb{R}\) satisfying
\[
\phi_L(\gamma_n) = 2\pi n,\qquad n=1,2,\dots
\]
Such a function exists by monotone Hermite interpolation. The \emph{generalised Riemann helix} is the space curve
\[
C_L(t) = \bigl(\cos\phi_L(t),\; \sin\phi_L(t),\; t\bigr),\qquad t\ge\gamma_1.
\]
As before, every zero lies on the generator \((1,0,t)\), and the curve winds exactly once between successive zeros.

The torsion of \(C_L\) is given by the Frenet--Serret formula
\[
\tau_L(t) = \frac{\phi_L'[(\phi_L')^4 - \phi_L'\phi_L''' + 3(\phi_L'')^2]}{(\phi_L')^6 + (\phi_L')^4 + (\phi_L'')^2}.
\tag{2}
\]

\section{Torsion from the Explicit Formula}

To connect with number theory we need the oscillatory part of \(\phi_L'\). From the zero‑alignment condition and the fact that the mean density of zeros is governed by the logarithmic derivative of the Gamma factors, we obtain
\[
\phi_L'(t) = \frac12 \sum_{j=1}^d \log\frac{t}{2\pi} + \text{const} + \text{osc}(t),
\]
where the oscillatory part \(\text{osc}(t)\) is obtained by differentiating the smooth interpolation of the zero‑counting function. The Weil explicit formula~\eqref{eq:weil} translates the distributional identity
\[
\sum_n \delta(t-\gamma_n) = \text{smooth} + \sum_{p^m} \frac{c(p^m)}{p^{m/2}} \delta(t - m\log p)
\]
into a pointwise statement for the regularised density. Concretely, introducing a narrow bump \(\eta_\varepsilon\), we have
\[
\phi_{L,\mathrm{osc}}'(t) = -2\sum_{p^m} \frac{c(p^m)}{p^{m/2}}\,\eta_\varepsilon(t - m\log p) + O(\varepsilon).
\tag{3}
\]
Higher derivatives \(\phi_L'', \phi_L'''\) involve \(\eta'\) and \(\eta''\). Substituting these into the torsion formula~\eqref{eq:torsion} and letting \(\varepsilon\to0\), the torsion degenerates into a distribution supported on the arithmetic points:
\[
\tau_L(t) = \sum_{p^m} \frac{c(p^m)}{p^{m/2}}\,\delta(t - m\log p) + \text{smooth background}.
\tag{4}
\]
Thus the torsion of the generalised helix carries the complete arithmetic information of the \(L\)-function.

\section{Logarithmic Compactification and Dirac Operator}

Change variable to \(x = \log t\), so that the peaks are located at \(x = m\log p\). Impose the periodicity \(x \sim x + L\) with \(L = 2\log\lambda\), where \(\lambda\) is a large scaling parameter (which one may take to be the analytic conductor of \(L(s)\)). On the resulting circle \(\mathbb{T}_L = \R/L\Z\) only finitely many prime powers \(p^m \le \lambda^2\) appear. The quantised geodesic flow on the unit tangent bundle of the compactified cylinder produces a Dirac operator with point interactions:
\[
D_{L,\lambda} = i\frac{d}{dx} + \sum_{p^m\le \lambda^2} \alpha_{p^m}\,\delta(x - m\log p \bmod L),
\qquad \alpha_{p^m} = \frac{c(p^m)}{p^{m/2}}.
\tag{5}
\]
The matching conditions at each singular point are
\[
\psi(x_j^+) = e^{i\alpha_{p^m}}\,\psi(x_j^-),
\]
with appropriate adjustments for the derivative. By the theory of self‑adjoint extensions for Dirac operators with point interactions~\cite{Albeverio1988}, \(D_{L,\lambda}\) is essentially self‑adjoint on a natural domain and has purely discrete spectrum.

\section{Spectral Determinant and Trace Formula}

Let \(\Delta_\lambda(z)\) be the regularised spectral determinant of \(D_{L,\lambda}\). Using the Krein resolvent formula, one can relate \(\Delta_\lambda(z)\) to the finite‑dimensional matrix of the truncated Weil quadratic form, as in the original CCM construction for \(\zeta(s)\). The generalised conjecture is:
\[
\lim_{\lambda\to\infty} \Delta_\lambda(z)\;\cdot\;(\text{explicit archimedean factor}) \;=\; \Lambda(z),
\tag{6}
\]
where \(\Lambda(z)\) is the completed \(L\)-function. If this holds, the spectral action \(\operatorname{Tr}(h(D_{L,\infty}))\) can be evaluated by a contour integral:
\[
\operatorname{Tr}(h(D_{L,\infty})) = \frac{1}{2\pi i} \oint h(z)\,\frac{\Lambda'(z)}{\Lambda(z)}\,dz + \text{trivial zeros}.
\]
The right‑hand side exactly equals the zero sum plus archimedean terms appearing in the explicit formula~\eqref{eq:weil}. Because the Weil class of test functions is a determining class, the spectrum of \(D_{L,\infty}\) must coincide with the set \(\{\pm\gamma_n\}\) of imaginary parts of the zeros of \(L(s)\).

\section{Dictionary}

\begin{center}
\begin{tabular}{c|c|c}
\textbf{Object} & \textbf{Riemann \(\zeta\)} & \textbf{General \(L\)-function} \\
\hline
Zeros & \(\gamma_n\) of \(\zeta(s)\) & \(\gamma_n\) of \(L(s)\) \\
Phase function & \(\phi(\gamma_n)=2\pi n\) & \(\phi_L(\gamma_n)=2\pi n\) \\
Arithmetic weights & \(\dfrac{\log p}{\sqrt{p}}\) & \(\dfrac{c(p^m)}{p^{m/2}}\) \\
Torsion peaks at & \(\log p\) & \(m\log p\) \\
Dirac coupling \(\alpha\) & \(\dfrac{\log p}{\sqrt{p}}\) & \(\dfrac{c(p^m)}{p^{m/2}}\) \\
Determinant limit & \(\Xi(z)\) (completed \(\zeta\)) & \(\Lambda(z)\) (completed \(L\)) \\
Spectrum & \(\pm\gamma_n\) & \(\pm\gamma_n\) \\
\end{tabular}
\end{center}

\section{Conclusion}

The compactified Riemann helix is the degree‑1 case of a universal geometric construction that pairs any automorphic \(L\)-function with a Dirac operator whose point‑interaction potential is dictated by the explicit formula. The eigenvalues of this operator are conjecturally the imaginary parts of the zeros, providing a uniform Hilbert–P\'olya realisation. While the core analytic difficulty (proving the convergence of the spectral determinants) remains open and is as deep as the Riemann Hypothesis for each \(L\)-function, the geometric blueprint unifies the problem and suggests a clear route for future work in spectral geometry and non‑commutative number theory.

\begin{thebibliography}{9}
\bibitem{CCM2025}
A.~Connes, C.~Consani, H.~Moscovici,
\emph{The scaling operator and a spectral triple for the Riemann zeta function},
arXiv:2511.22755 (2025).

\bibitem{Albeverio1988}
S.~Albeverio, F.~Gesztesy, R.~H{\o}egh-Krohn, H.~Holden,
\emph{Solvable Models in Quantum Mechanics},
Springer, 1988.

\bibitem{Weil1952}
A.~Weil,
\emph{Sur les formules explicites de la th\'eorie des nombres},
Comm. Sem. Math. Univ. Lund, tome suppl. d\'edi\'e \`a Marcel Riesz (1952), 252--265.

\bibitem{Selberg1956}
A.~Selberg,
\emph{Harmonic analysis and discontinuous groups in weakly symmetric Riemannian spaces with applications to Dirichlet series},
J.~Indian Math.~Soc.~\textbf{20} (1956), 47--87.
\end{thebibliography}

\end{document}